% --- Kurzbiografien der Autor*innen ---
% Sortierung: alphabetisch nach Nachname.
\backmattersection{Kurzbiografien}

Im Folgenden finden Sie kurze biografische Angaben zu allen Autor*innen, deren Beiträge in diesem Band versammelt sind. Die Reihenfolge ist alphabetisch nach Nachname.

\subsection*{Anke Albrecht}
\noindent\small Friedrich-Alexander-Universität Erlangen-Nürnberg \\
Anke Albrecht ist wissenschaftliche Mitarbeiterin am Lehrstuhl für Digital Humanities der FAU Erlangen-Nürnberg und arbeitet seit über sechs Jahren mit WissKI in unterschiedlichen Forschungskontexten. Ihre Schwerpunkte liegen in der semantischen Modellierung historischer Korrespondenznetze und in der Anbindung an internationale Normdatenanbieter. Aktuell forscht sie zur ontologischen Repräsentation textueller Werkfassungen und ist Mitautorin mehrerer Konfigurationsempfehlungen, die im WissKI-Konsortium für neu startende Projekte abgestimmt werden.

\subsection*{Bernd Braun}
\noindent\small Friedrich-Alexander-Universität Erlangen-Nürnberg \\
Bernd Braun ist technischer Koordinator des WissKI-Konsortiums an der FAU Erlangen-Nürnberg. Seine Schwerpunkte liegen in der Weiterentwicklung des WissKI-Pathbuilders, in der Konfiguration und Inbetriebnahme produktiver Installationen sowie in der Unterstützung neu startender Projekte. Er war an der Entwicklung der aktuellen WissKI-Konfigurationsvorlage maßgeblich beteiligt und engagiert sich in der Open-Source-Pflege des WissKI-Quellcodes.

\subsection*{Clara Castorp}
\noindent\small Universität Heidelberg \\
Clara Castorp ist Wissenschaftlerin am Heidelberger Zentrum für Digital Humanities und beschäftigt sich seit zehn Jahren mit der Migration historisch gewachsener Sammlungsdatenbanken in semantische Forschungsumgebungen. Sie hat zahlreiche Migrationsprojekte für Sammlungseinrichtungen unterschiedlicher Größenordnung wissenschaftlich begleitet und ist Mitautorin einer einschlägigen Methodendokumentation zur Datenbereinigung in der Sammlungspraxis. Aktuell arbeitet sie an einem Forschungsdatenmanagement-Konzept für historische Wissenschaftssammlungen.

\subsection*{Daniel Drechsler}
\noindent\small Humboldt-Universität zu Berlin \\
Daniel Drechsler ist Informatiker mit einem Hintergrund in der Wissensrepräsentation und arbeitet als Forschungsdatenmanager an der Humboldt-Universität zu Berlin. Sein Aufgabenfeld umfasst die Begleitung geisteswissenschaftlicher Forschungsprojekte beim Übergang in semantische Datenmodelle. Er entwickelt seit mehreren Jahren ETL-Werkzeuge für die Triple-Store-Migration und hat dazu mehrfach in einschlägigen Foren der digitalen Geisteswissenschaften vorgetragen.

\subsection*{Elena Engel}
\noindent\small Universität Wien \\
Elena Engel ist wissenschaftliche Mitarbeiterin am Institut für Kunstgeschichte der Universität Wien. Ihre Forschungsschwerpunkte liegen im Bereich der digitalen Bildwissenschaft, insbesondere in der Integration von IIIF-basierten Bildrepositorien in semantische Forschungsumgebungen. Sie ist beteiligt an mehreren europäischen Forschungsprojekten zur Interoperabilität kunsthistorischer Bilddatenbestände und engagiert sich in der IIIF-Community in der Arbeitsgruppe \glqq Cultural Heritage\grqq.

\subsection*{Frank Förster}
\noindent\small Universität zu Köln \\
Frank Förster ist Bibliothekar an der Universitäts- und Stadtbibliothek Köln und betreut dort den Bereich Digital Humanities. Seine Schwerpunkte liegen in der Lizenzierung digitaler Bildressourcen, in der bibliotheksnahen Anwendung von IIIF und in der semantischen Erschließung historischer Bestände. Er ist Mitherausgeber einer Themenreihe zu Rechtsfragen der digitalen Bildpublikation und engagiert sich in der bibliothekarischen Standardisierungsarbeit.

\subsection*{Greta Gärtner}
\noindent\small Friedrich-Alexander-Universität Erlangen-Nürnberg \\
Greta Gärtner ist Postdoktorandin an der FAU Erlangen-Nürnberg und Mitglied des WissKI-Konsortiums. Ihre Forschungsschwerpunkte liegen in der Anwendung semantischer Datenmodelle auf museale Sammlungen mit kleinem und mittelgroßem Bestand. Sie hat in mehreren Migrationsprojekten als wissenschaftliche Begleitung mitgewirkt und beschäftigt sich aktuell mit der Schnittstellenanbindung von WissKI-Beständen an Museums-Verbund-Plattformen.

\subsection*{Hannes Hauser}
\noindent\small Technische Universität Dresden \\
Hannes Hauser ist Forschungsdatenmanager am Sächsischen Landesdigitalisierungszentrum an der TU Dresden. Er betreut SPARQL-Endpunkte und Datenpublikations-Workflows für mehrere geisteswissenschaftliche Forschungsprojekte und entwickelt im Rahmen einer Promotionsarbeit ein Monitoring-Framework für die Langzeitstabilität semantischer Forschungsdatendienste.

\subsection*{Iris Ihlow}
\noindent\small Universität Hamburg \\
Iris Ihlow ist Kunsthistorikerin und arbeitet als Sammlungsleitung an einem mittelgroßen kulturhistorischen Museum. Über eine Honorarprofessur ist sie an die Universität Hamburg angebunden und vermittelt dort museale Sammlungsforschung mit digitalen Methoden. Sie verantwortet seit fünf Jahren die LOD-Strategie ihres Hauses und engagiert sich in der Arbeitsgruppe Digitalisierung des Deutschen Museumsbundes.

\subsection*{Jan Jansen}
\noindent\small Rheinische Friedrich-Wilhelms-Universität Bonn \\
Jan Jansen ist Wissenschaftlicher Mitarbeiter im Bereich Forschungsdatenmanagement an der Universität Bonn. Sein Aufgabenfeld umfasst die Beratung geisteswissenschaftlicher Forschungsprojekte bei der Konzeption von Datenmanagementplänen und die institutionelle Verstetigung von LOD-Diensten. Er engagiert sich in der politischen Debatte um die Finanzierung digitaler Forschungsinfrastrukturen.

\subsection*{Karin Krüger}
\noindent\small Georg-August-Universität Göttingen \\
Karin Krüger ist Editionswissenschaftlerin an der Universität Göttingen und leitet seit zwei Förderperioden ein editionsphilologisches Forschungsprojekt zu Gelehrtenkorrespondenzen des 18.\,Jahrhunderts. Ihre Schwerpunkte liegen in der Anbindung editionswissenschaftlicher Forschungsdaten an Normdatenanbieter und in der methodischen Reflexion über das Verhältnis von Editionspraxis und semantischer Datenmodellierung.

\subsection*{Lukas Lehmann}
\noindent\small Berlin-Brandenburgische Akademie der Wissenschaften \\
Lukas Lehmann ist Softwareentwickler an der Berlin-Brandenburgischen Akademie der Wissenschaften und betreut dort mehrere WissKI-basierte Editionsvorhaben. Sein Aufgabenfeld umfasst die Entwicklung von WissKI-Erweiterungen, die Anbindung an externe Normdaten- und Recherchedienste und die laufende Wartung der produktiven Forschungsdatensysteme der Akademie.

\subsection*{Mara Möller}
\noindent\small Eberhard Karls Universität Tübingen \\
Mara Möller ist Wissenschaftlerin am Tübinger Zentrum für Digital Humanities. Sie beschäftigt sich seit über zehn Jahren mit Versionierungs- und Archivierungsfragen in semantischen Forschungsdatenkontexten und hat mehrere institutionelle Workflows zur Zitierfähigkeit von WissKI-Datenbeständen entwickelt. Aktuell forscht sie zu Standards der Langzeitarchivierung graphbasierter Forschungsdaten.

\subsection*{Nadja Neumann}
\noindent\small Universitäts- und Landesbibliothek Münster \\
Nadja Neumann ist Leiterin der Abteilung Forschungsdatendienste an der Universitäts- und Landesbibliothek Münster. Sie verantwortet die institutionelle Langzeitarchivierung geisteswissenschaftlicher Forschungsdaten und betreut die Übergabe von Forschungsprojekten in den Daueretat der Bibliothek. Ihr Engagement in der bibliothekarischen Standardisierungsarbeit umfasst mehrere Arbeitsgruppen der DINI und der DFG.

\subsection*{Olaf Otten}
\noindent\small Friedrich-Alexander-Universität Erlangen-Nürnberg \\
Olaf Otten ist Lehrbeauftragter am Lehrstuhl für Digital Humanities der FAU Erlangen-Nürnberg und Autor mehrerer Lehrbücher zur Ontologiemodellierung. Er verantwortet seit acht Jahren das Lehrprogramm des WissKI-Konsortiums zur Aus- und Weiterbildung in semantischen Datenmodellen und gibt regelmäßig Workshops zu Protégé und WissKI.

\subsection*{Petra Pohl}
\noindent\small Universität Leipzig \\
Petra Pohl ist Systemadministratorin und IT-Beraterin am Universitätsrechenzentrum Leipzig. Sie betreut die produktiven Drupal- und WissKI-Installationen der Leipziger Geisteswissenschaften und hat in dieser Funktion zahlreiche Neuinstallationen begleitet. Ihre Schwerpunkte liegen in der Inbetriebnahme nachhaltiger Server-Konfigurationen und in der Vermittlung des dafür benötigten Wissens.

\subsection*{Rita Reuter}
\noindent\small Deutsches Literaturarchiv Marbach \\
Rita Reuter ist Mitarbeiterin der Abteilung Digital Humanities am Deutschen Literaturarchiv Marbach. Sie betreut SPARQL-basierte Recherche-Werkzeuge der Marbacher Forschungsdaten und entwickelt didaktische Konzepte zur Vermittlung von Abfragesprachen für Geisteswissenschaftler*innen. Sie ist Autorin mehrerer einschlägiger Aufsätze zur Didaktik der digitalen Geisteswissenschaften.

\subsection*{Stefan Schubert}
\noindent\small Johannes Gutenberg-Universität Mainz \\
Stefan Schubert ist Wissenschaftler an der Johannes Gutenberg-Universität Mainz und Autor eines im erscheinenden Lehrbuchs zur Abfragesprache SPARQL für die Geisteswissenschaften. Er hält regelmäßig Lehrveranstaltungen zu semantischen Datenmodellen und engagiert sich in der Entwicklung offener Lehrmaterialien für die digitalen Geisteswissenschaften.

\subsection*{Tanja Thiel}
\noindent\small Julius-Maximilians-Universität Würzburg \\
Tanja Thiel ist Wissenschaftliche Mitarbeiterin am Institut für Sammlungswissenschaft der Universität Würzburg. Ihre Schwerpunkte liegen in der ereignisorientierten Sammlungsmodellierung nach CIDOC CRM und in der didaktischen Vermittlung dieser Modellierungsperspektive an Sammlungspraktiker*innen. Sie ist Autorin einer einschlägigen Einführung in das CIDOC CRM für die Sammlungspraxis.

\subsection*{Udo Ullmann}
\noindent\small Albert-Ludwigs-Universität Freiburg \\
Udo Ullmann ist Wissenschaftler am Lehrstuhl für Informationsvisualisierung der Universität Freiburg. Seine Forschungsschwerpunkte liegen in der Visualisierung großer geisteswissenschaftlicher Datenbestände und in der methodischen Reflexion über die argumentative Funktion von Visualisierungen. Er ist Autor mehrerer Bücher zur Visualisierungspraxis in den digitalen Geisteswissenschaften.

\subsection*{Vera Vogt}
\noindent\small Friedrich-Alexander-Universität Erlangen-Nürnberg \\
Vera Vogt ist wissenschaftliche Mitarbeiterin am Lehrstuhl für Digital Humanities der FAU Erlangen-Nürnberg. Ihre Schwerpunkte liegen in der Anwendung von Graphen-Visualisierungsverfahren auf semantische Forschungsdatenbestände und in der Werkzeugauswahl für die digitale Forschungspraxis. Sie engagiert sich in der Aus- und Weiterbildung des WissKI-Konsortiums.

\subsection*{Wolfgang Wagner}
\noindent\small Universität Stuttgart \\
Wolfgang Wagner ist Germanist und Editionsphilologe an der Universität Stuttgart. Er leitet seit zwei Jahren ein DFG-gefördertes Editionsvorhaben zur Korrespondenz Wilhelm Müllers und beschäftigt sich darüber hinaus mit Fragen der Plattformwahl in digitalen Editionsprojekten. Sein Engagement umfasst mehrere Beiräte im Umfeld der digitalen Editionsphilologie.
